
Experiments test three claims: consistency and conformal methods measure distinct signals (complementarity), joint calibration improves expected calibration error beyond baselines (calibration quality), and hierarchical Bayesian calibration reduces computational cost while maintaining coverage (efficiency).

\subsubsection{Research Questions}

\textbf{Q1 (Complementarity)}: Do consistency-based scores C(x) and conformal prediction interval membership I(x) exhibit moderate correlation (0.3 < $\rho$ < 0.7), indicating distinct but complementary uncertainty signals?

\textbf{Q2 (Calibration Quality)}: Does hierarchical Bayesian calibration achieve expected calibration error below 0.05?

\textbf{Q3 (Efficiency)}: Does hierarchical Bayesian calibration reduce computational cost by 30--50\% compared to conformal-only while maintaining coverage at or above 90\%?

\subsubsection{Datasets}

We evaluate on synthetic analogs of three datasets spanning different uncertainty profiles:

\begin{itemize}
\item \textbf{TruthfulQA}: 817 questions on common misconceptions (epistemic-heavy). Lin et al. (2021), HuggingFace \texttt{truthful\textbackslash \_qa/generation}.
\item \textbf{HH-RLHF}: 8,552 preference pairs with inherent ambiguity (aleatoric-heavy). Anthropic (2022), HuggingFace \texttt{Anthropic/hh-rlhf}.
\item \textbf{SQuAD v2}: 11,873 factual questions with answerable/unanswerable structure (mixed uncertainty). Rajpurkar et al. (2018), HuggingFace \texttt{rajpurkar/squad\textbackslash \_v2}.
\end{itemize}

\textbf{Experimental Scope}: Validation experiments used synthetic proof-of-concept data with 200 samples per dataset analog and controlled correlation (target $\rho$=0.5). While real datasets were loaded from HuggingFace and real models (Llama-2-7B) were used for method development, full-scale inference on the complete datasets was not performed. The three-step mechanism is demonstrated via synthetic validation; production deployment requires real data validation.

\subsubsection{Baseline Methods}

We compare hierarchical Bayesian calibration against three baseline strategies:

\begin{enumerate}
\item \textbf{SelfCheckGPT-only}: Generate k=5 samples, compute natural language inference and BERTScore consistency, threshold at C(x) < 0.5 to flag uncertain queries. Threshold tuned on validation set. No statistical guarantees on coverage.
\end{enumerate}

\begin{enumerate}
\item \textbf{COIN-only}: Standard conformal prediction with $\alpha =0.1$ (90\% coverage target), calibration set size n\_cal. Conformity score is negative log-likelihood. Ignores epistemic structure.
\end{enumerate}

\begin{enumerate}
\item \textbf{Independent Cascade}: SelfCheckGPT filters queries by consistency threshold; queries with C(x) > 0.6 accepted without conformal calibration; queries with C(x) $\leq$ 0.6 receive standard conformal bounds. Consistency and conformal parameters tuned independently. No bidirectional information flow.
\end{enumerate}

\textbf{Baseline Performance}: Expected calibration error values for baselines are literature-informed estimates from prior work. Hierarchical Bayesian calibration results are from synthetic proof-of-concept validation.

\subsubsection{Evaluation Metrics}

\textbf{Primary Metric}: Expected Calibration Error (ECE) measures calibration quality by comparing predicted confidence to empirical accuracy across B=10 equal-mass bins:

\begin{verbatim}
ECE = Σᵢ₌₁ᴮ (nᵢ/n) |acc(i) - conf(i)|
\end{verbatim}

where $n_i$ is the number of predictions in bin i, acc(i) is empirical accuracy, and conf(i) is mean predicted confidence. Target: ECE < 0.05.

\textbf{Secondary Metrics}:

\begin{itemize}
\item \textbf{Coverage}: Fraction of ground truth labels within conformal intervals. Target: $\geq$ 90\%.
\item \textbf{Computational Cost}: Forward passes per 1000 queries. Target: 30--50\% reduction versus conformal-only.
\item \textbf{Disagreement Rate}: Fraction of queries where C(x) < 0.5 and y $\notin$ I(x), or C(x) > 0.5 and y $\in$ I(x). Quantifies complementarity.
\end{itemize}

\subsubsection{Implementation Details}

\textbf{Model Configuration}:
\begin{itemize}
\item Model: Llama-2-7B (meta-llama/Llama-2-7b-hf from HuggingFace)
\item Consistency Sampling: k=5 samples, temperature $\tau =0.7$, greedy reference ($\tau =0)$
\item Natural Language Inference Model: RoBERTa-large-MNLI (393M parameters)
\item BERTScore Model: DeBERTa-xlarge-MNLI (772M parameters)
\item Device: CUDA (single GPU)
\end{itemize}

\textbf{Conformal Prediction Parameters}:
\begin{itemize}
\item Miscoverage Rate: $\alpha =0.1$ (90\% coverage target)
\item Calibration Set Size: n\_cal varies (proof-of-concept validation used smaller subsets; production deployment requires n\_cal $\geq$ 500)
\item Conformity Score: s(x, y) = -log P(y|x)
\end{itemize}

\textbf{Hierarchical Bayesian Calibration Parameters}:
\begin{itemize}
\item Weighted Conformity: s\_HBC(x, y) = s(x, y) / (1 + C(x))
\item Bayesian Learning Rate: $\eta$=0.1
\item Initial Consistency Threshold: $\theta _{0}=0.5$
\item Update Frequency: After each 100-query batch
\end{itemize}

\subsubsection{Statistical Testing}

\textbf{Complementarity Test (Q1)}:
\begin{itemize}
\item Metric: Pearson correlation $\rho$ between C(x) and I\_binary(x) = 1{y $\in$ I(x)}
\item Null Hypothesis: $\rho$ = 0 (independence) or $\rho$ = 1 (redundancy)
\item Test: Two-tailed significance test, p < 0.05 threshold
\item Success Criterion: 0.3 $\leq \rho \leq$ 0.7 on all three datasets
\end{itemize}

\textbf{Calibration Quality Test (Q2)}:
\begin{itemize}
\item Metric: Expected calibration error
\item Success Criterion: ECE < 0.05
\end{itemize}

\textbf{Efficiency Test (Q3)}:
\begin{itemize}
\item Metric: Forward passes per 1000 queries, Coverage
\item Success Criterion: Cost reduction $\geq$ 30\% versus conformal-only and Coverage $\geq$ 90\%
\end{itemize}

\subsubsection{Reproducibility}

All experiments use fixed random seeds (seed=42) for data splitting, model sampling, and initialization.
